
Quality-first curation shows universal superiority across tested scales (10K--10M tokens, Cohen's d=0.76--2.48), providing evidence for quality saturation mechanisms and suggesting quality-gated diversity effects. We quantify quality saturation (slope $-2.62pp/log-scale$, p=0.008) and model diversity persistence based on FAC literature (slope +2.5pp/log-scale, simulated). More critically, we show evidence for compositional ordering effects suggesting that sequential curation creates non-additive interactions, extending Data Mixing Laws' domain proportion framework to sequential filtering steps.

The core insight is quality-gated diversity: diversity sampling appears most effective on quality-filtered subsets. When diversity precedes quality (DQ), diverse sampling operates on full noisy distribution, potentially amplifying uninformative variation. When quality precedes diversity (QD), filtering removes noise first, allowing diversity selection to preserve genuinely informative patterns.

For practitioners working with 100M--1B parameter models on web corpora at 10K--10M scale: apply quality filtering before diversity sampling. This recommendation holds even at large scale where diversity coefficient exceeds quality coefficient, because compositional dependency appears to favor quality-first ordering.

Future work includes extended scale tests (50M--100M tokens) to search for reversal threshold, real diversity experiments (h-e2 with full FAC implementation) to validate simulated trajectory, and frontier model replication (Llama 3 8B) to validate proxy transfer assumptions.

Sequential ordering appears to matter in data curation---apply quality filtering before diversity sampling, not because diversity is unimportant, but because diversity appears most effective when noise is removed first.
